\documentclass[10pt]{article}

\usepackage[margin=1in]{geometry}
\usepackage{times}
% microtype removed: causes \showhyphens conflict with times in some distributions
\usepackage{amsmath,amssymb,amsthm}
\usepackage{algorithm}
\usepackage{algorithmic}
\usepackage{enumitem}
\usepackage{graphicx}
\usepackage[hidelinks]{hyperref}
\usepackage[numbers]{natbib}
\usepackage{booktabs}
\usepackage{multirow}
\usepackage[table]{xcolor}
\usepackage{listings}
\usepackage{caption}

\lstset{
  basicstyle=\small\ttfamily,
  breaklines=true,
  frame=single,
  numbers=left,
  numberstyle=\tiny,
  tabsize=2
}

\theoremstyle{definition}
\newtheorem{definition}{Definition}[section]

\title{Interactive Constraint Explanations for Human--Agent Graph Colouring:\\
       A Comparative Study of Perspective and Representation}
\author{}
\date{}

\begin{document}
\maketitle

\begin{abstract}
Effective collaboration between humans and AI systems in constraint-satisfaction tasks depends not
only on finding valid solutions, but on helping human collaborators understand the constraint
structure that governs feasibility.  We present a platform for studying how different
representations of constraint information affect human understanding, workload, and task performance
in a graph-colouring task with partial observability.  Rather than requiring the human to negotiate
with agents via dialogue, the redesigned interface exposes the implications of the human's choices
in real time: agents act as \emph{constraint analysers} that continuously compute and display the
feasible solution space induced by the human's current node-colour assignments.  Two design
dimensions define four experimental conditions: \emph{perspective} (user-centric vs.\
agent-centric framing of constraint information) and \emph{representation} (formulaic logical
expressions vs.\ natural-language summaries generated by an LLM).  We describe the formal model,
the real-time computation engine, the four interface conditions in detail, and a user-study design
comparing objective convergence metrics against subjective workload, clarity, and trust.
\end{abstract}

%% ============================================================
\section{Introduction}
\label{sec:intro}
%% ============================================================

Humans working alongside AI agents in constraint-satisfaction tasks face a fundamental
interpretive challenge: understanding how their own decisions constrain the options available to
their collaborators.  In classical approaches, agents negotiate solutions with human partners
through structured or conversational dialogue.  While such systems have been studied extensively,
dialogue introduces its own cognitive demands --- interpreting agent moves, formulating responses,
tracking negotiation state --- which are orthogonal to the underlying optimisation problem.

This paper re-centres the research question from \emph{how should agents negotiate with humans?}
to \emph{how should agents explain constraint structure to humans?}  We hypothesise that the
principal difficulty in human--agent coordination is not reaching agreement through dialogue but
helping humans understand the constraint space in which they are operating.

We investigate this question using a graph-colouring task in which the human directly manipulates
node colours and the system continuously reveals the implications of those choices.  Agents no
longer act as negotiating partners; instead, they act as constraint analysers.  As the human
adjusts their cluster's colours, the system recomputes --- in real time --- which configurations
remain feasible for each agent cluster, and displays this information through one of four
interface conditions.

The two independent variables are:
\begin{enumerate}
  \item \textbf{Perspective}: whether constraint information is framed relative to the
        \emph{human's choices} (user-centric) or relative to the \emph{agent's remaining options}
        (agent-centric).
  \item \textbf{Representation}: whether constraints are expressed as \emph{formulaic logical
        statements} or as \emph{natural-language summaries} produced by an LLM.
\end{enumerate}

The contributions of this paper are:
\begin{enumerate}
  \item A formal model of interactive constraint visualisation for distributed graph colouring,
        including definitions of feasible sets, domain projections, and infeasibility conditions.
  \item A real-time constraint computation engine that exactly evaluates agent feasibility given
        arbitrary partial human assignments.
  \item Detailed specifications of four interface conditions spanning two design dimensions.
  \item A user-study design for comparing these conditions on objective and subjective measures.
\end{enumerate}

%% ============================================================
\section{Related Work}
\label{sec:related}
%% ============================================================

\subsection{Distributed Constraint Optimisation}
Distributed Constraint Optimisation Problems~\citep{modi2005adopt} provide a canonical framework
for multi-agent coordination under coupled local constraints.  Classic algorithms such as
DPOP~\citep{petcu2005dpop} and Max-Sum~\citep{rogers2011bounded} achieve globally optimal or
bounded-suboptimal solutions through structured message passing.  Our work adopts the DCOP
formulation as the experimental task but shifts the focus from algorithm design to human
understanding of the feasibility structure that such algorithms exploit.

\subsection{Mixed-Initiative and Human--AI Collaboration}
Mixed-initiative systems allow control to be shared adaptively between humans and
agents~\citep{horvitz1999principles}.  Within optimisation, this raises the question of how much
constraint information the human needs to make useful contributions.  Approaches range from
fully automated solving with post-hoc explanation to interactive solving where humans manipulate
partial solutions and agents fill in the rest~\citep{green2019humanintheloop}.  Our interface
occupies the latter end: the human has full control over their cluster and agents respond by
exposing the downstream implications rather than issuing demands.

\subsection{Explanations in AI Systems}
The explainability literature distinguishes \emph{global} explanations (properties of the model
as a whole) from \emph{local} explanations (reasons for a specific output)~\citep{ribeiro2016why}.
Closest to our setting is contrastive explanation~\citep{miller2019explanation}, which answers the
question ``why this outcome rather than another?''  User-centric constraint framing is a form of
contrastive explanation: ``because you chose $h_4 = \textit{blue}$, the agent is forced into a
specific subset of configurations.''  We operationalise and compare such contrastive framing
against a more direct display of remaining options (agent-centric).

\subsection{Constraint Visualisation}
Several interactive constraint-solving tools have been developed to support human understanding of
feasibility~\citep{borning1997constraint,freuder2003progress}.  These systems typically expose
constraint propagation results visually, allowing users to observe which variable domains contract
as they fix values.  Our system extends this approach to a distributed multi-agent setting where
one agent's feasible set depends on another participant's choices, and where the explanation
interface is itself the experimental variable.

\subsection{LLMs for Explanation and Summarisation}
LLMs have demonstrated strong capability in summarising complex technical content in accessible
natural language~\citep{brown2020language}.  Recent work has applied LLMs to explain planning
and optimisation outputs to lay audiences~\citep{guo2025castl}.  A key design question is
whether LLM-generated summaries preserve sufficient precision for constraint reasoning, or whether
formulaic expressions are required for accurate interpretation.  Our experiment directly compares
these two representation styles under controlled conditions.

%% ============================================================
\section{Problem Formulation}
\label{sec:problem}
%% ============================================================

\subsection{The Graph-Colouring Task}

We use the same distributed graph-colouring task as the original study.  The global graph
$G = (V, E)$ is partitioned among three participants --- two artificial agents $a_1, a_2$ and one
human $h$ --- each controlling a disjoint cluster of nodes:
$V = V_{a_1} \sqcup V_h \sqcup V_{a_2}$.
Each participant assigns a colour from a shared domain $\mathcal{C}$ (default: red, green, blue)
to each of their nodes.  The global conflict penalty is:
\begin{equation}
  J(x) = \sum_{(u,v) \in E} \mathbf{1}[x_u = x_v]
  \label{eq:penalty}
\end{equation}
where $x: V \to \mathcal{C}$ is the global assignment.

In the default experiment: $|V| = 15$ nodes (five per cluster), $|\mathcal{C}| = 3$,
approximately 4--6 inter-cluster boundary edges, and a cluster topology of
$V_{a_1} \leftrightarrow V_h \leftrightarrow V_{a_2}$ (the two agent clusters are coupled only
through the human's cluster).

\subsection{Partial Observability}

Each participant has \emph{local} visibility: they fully observe their own cluster and can see
the boundary-neighbour nodes that share an edge with their cluster, but have no knowledge of
nodes further inside other clusters.  Formally, participant $p$'s visible external nodes are:
\[
  \partial V_p = \{v \notin V_p \mid \exists\, u \in V_p: (u,v) \in E\}
\]
No topology beyond $\partial V_p$ is ever exposed to participant $p$, whether in the UI or in any
generated explanation.

\subsection{Partial Human Assignment}

In the new interface the human's configuration is modelled as a \emph{partial} assignment, since
nodes may be left unspecified (displayed as grey) while the human explores possibilities:
\begin{equation}
  \alpha_h : V_h \to \mathcal{C} \cup \{\bot\}
  \label{eq:partial}
\end{equation}
where $\bot$ represents an unassigned (grey) node.  The \emph{effective domain} of $\alpha_h$
is $\mathrm{dom}(\alpha_h) = \{v \in V_h \mid \alpha_h(v) \neq \bot\}$.

This differs from the original negotiation study, where the human always maintained a complete
assignment.  Allowing grey nodes is essential for exploration: the human can fix a subset of their
choices and observe feasibility before committing to the rest.

\subsection{Agent Feasibility Sets}

Given the human's partial assignment $\alpha_h$, the set of \emph{feasible configurations} for
agent cluster $a$ is:
\begin{equation}
  S_a(\alpha_h) = \bigl\{ x_a \in \mathcal{C}^{|V_a|}
    \;\big|\; J_a(x_a, \alpha_h) = 0 \bigr\}
  \label{eq:feasible}
\end{equation}
where $J_a(x_a, \alpha_h)$ counts conflicts on edges incident to $V_a$, treating unassigned
boundary-neighbour nodes as imposing no conflict constraint.  The system computes $S_a(\alpha_h)$
\emph{exactly} by exhaustive enumeration over $|\mathcal{C}|^{|V_a|}$ candidates (tractable for
$k=3$, $n=5$: 243 combinations).

The cluster is \emph{feasible} if $S_a(\alpha_h) \neq \emptyset$ and \emph{infeasible} if
$S_a(\alpha_h) = \emptyset$.  Infeasibility means no valid global colouring exists given the
human's current choices; the interface notifies the human and proposes a minimal repair
(Section~\ref{sec:feasibility}).

\subsection{Constraint Projections}

Because the full feasible set $S_a(\alpha_h)$ is a set of complete cluster assignments, it is
difficult to display directly.  The interface instead presents two \emph{projections}.

\paragraph{Domain projection (agent-centric).}
For each agent node $u \in V_a$, the set of colours that $u$ can still adopt:
\begin{equation}
  D_u(\alpha_h) = \bigl\{ c \in \mathcal{C} \;\big|\;
    \exists\, x_a \in S_a(\alpha_h): x_a(u) = c \bigr\}
  \label{eq:domain}
\end{equation}
This is the marginal projection of $S_a(\alpha_h)$ onto node $u$.

\paragraph{Consequence projection (user-centric).}
For each human node $v \in V_h$ whose colour is fixed, the subset of agent configurations
that are reachable:
\begin{equation}
  S_a(\alpha_h \mid v = c) = \bigl\{ x_a \in S_a(\alpha_h)
    \;\big|\; \text{consistent with } \alpha_h(v) = c \bigr\}
  \label{eq:consequence}
\end{equation}
The interface presents this as a logical description of which agent-node assignments are
jointly implied (or excluded) by the human's choice of $\alpha_h(v) = c$.

%% ============================================================
\section{System Architecture}
\label{sec:system}
%% ============================================================

\subsection{Graph Interface}

The participant sees a graph canvas displaying:
\begin{itemize}
  \item All nodes in the human cluster $V_h$, coloured according to $\alpha_h$ (grey if
        unassigned).
  \item The boundary-neighbour nodes $\partial V_h$ from each agent cluster, displayed but
        not directly manipulable by the human.
  \item Edges connecting nodes, with conflicting edges highlighted.
\end{itemize}

The human cycles through colours (red, green, blue, grey) by clicking any node in $V_h$.
Available colours are $\mathcal{C} = \{\text{red}, \text{green}, \text{blue}\}$; grey represents
the unassigned state $\bot$.

\subsection{Real-Time Constraint Engine}

Whenever the human changes a node colour, the system executes the following pipeline:
\begin{enumerate}
  \item Update $\alpha_h$ with the new assignment.
  \item For each agent $a \in \{a_1, a_2\}$, recompute $S_a(\alpha_h)$ by exhaustive enumeration.
  \item For each agent node $u$, recompute $D_u(\alpha_h)$ from $S_a(\alpha_h)$.
  \item Compute user-centric consequence descriptions for each fixed human node.
  \item Update all constraint panels with the new information.
\end{enumerate}
This pipeline completes within milliseconds for the default graph size, so updates appear
instantaneous.  There is no turn structure and no messaging interface.

\subsection{Constraint Panels}

The interface contains two types of constraint display panels, each populated according to the
active experimental condition:

\paragraph{Node-level panels.}
Each node (human or agent) may display an annotation showing its constraint status.  In
agent-centric conditions, agent nodes display their domain $D_u(\alpha_h)$.  In user-centric
conditions, human nodes display the consequence sets induced by their current colour choice.

\paragraph{Global feasibility panel.}
A summary panel at the top of the screen shows the feasibility status of each agent cluster.
Possible states are: \textsc{Feasible} ($S_a \neq \emptyset$), \textsc{Infeasible}
($S_a = \emptyset$), and \textsc{Unconstrained} (no boundary edges fixed yet, so all
configurations remain possible).

\subsection{Infeasibility Repair Proposals}
\label{sec:feasibility}

When $S_a(\alpha_h) = \emptyset$, the system identifies a minimal set of human nodes whose
reassignment would restore feasibility.  Formally, it finds the smallest set
$R \subseteq \mathrm{dom}(\alpha_h)$ such that:
\[
  S_a\bigl(\alpha_h \mid_{V_h \setminus R}\bigr) \neq \emptyset
\]
This is computed by iterating over subsets of fixed human nodes in order of increasing size.
The interface then proposes specific colour changes for the nodes in $R$, framed according to the
active condition (as a logical statement or as a natural-language suggestion).

\subsection{LLM Summary Generation}
\label{sec:llm}

In the LLM-representation conditions (Conditions 3 and 4), the formulaic constraint expressions
are passed to a language model to produce natural-language summaries.  The LLM receives:
\begin{itemize}
  \item The precomputed domain set $D_u(\alpha_h)$ or consequence description for the relevant node.
  \item The current human assignment $\alpha_h$.
  \item A system prompt instructing it to summarise the constraint information in accessible
        natural language without inventing additional constraints.
\end{itemize}
The LLM \emph{does not perform any reasoning or constraint computation}.  It acts as a
\emph{presentation layer only}: the constraint information it receives is exact and
precomputed; its task is solely to render it more readable.  All LLM calls are logged with the
input constraint data and output text, enabling post-hoc verification that no hallucinated
constraints were introduced.

%% ============================================================
\section{Experimental Conditions}
\label{sec:conditions}
%% ============================================================

The study uses a $2 \times 2$ factorial design with two within-subjects factors:

\begin{description}[leftmargin=2em]
  \item[Perspective] \emph{User-centric} (constraints framed relative to the human's choices)
        vs.\ \emph{Agent-centric} (constraints framed relative to agent nodes' remaining options).
  \item[Representation] \emph{Formulaic} (logical expressions) vs.\ \emph{Natural language}
        (LLM-generated summaries of the same information).
\end{description}

This produces four conditions.  All conditions share the same real-time computation engine and
the same graph interface; only the constraint panels differ.

\subsection{Condition~1: User-Centric Formulaic}
\label{sec:c1}

Human nodes display logical expressions describing the set of agent configurations that are
jointly induced by the human's current colour choices.

\paragraph{Display format.}
Each fixed human node $v$ with $\alpha_h(v) = c$ displays the constraint:
\begin{center}
\small
\texttt{Because: $v$ = $c$}\\
\texttt{Requires: $(\text{Agent configs consistent with } \alpha_h)$}
\end{center}
When $\alpha_h(v) = c$ forces a specific joint agent assignment, this is displayed as a
conjunction.  When multiple compatible agent configurations remain, they are listed as a
disjunction:
\begin{center}
\small
\texttt{Because: h4 = Blue}\\
\texttt{Means: (a5 = Red AND a4 = Green) OR (a5 = Green AND a4 = Red)}
\end{center}

\paragraph{Design intent.}
This condition provides the most informationally precise display, directly linking the human's
decisions to agent outcomes.  It is expected to be cognitively demanding, particularly when many
agent configurations remain compatible.

\subsection{Condition~2: Agent-Centric Formulaic}
\label{sec:c2}

Agent nodes display the set of colours they may still adopt, expressed as a logical formula.

\paragraph{Display format.}
Each agent node $u$ displays its domain projection $D_u(\alpha_h)$:
\begin{center}
\small
\texttt{a5: \{Red, Green\}}
\end{center}
When a node's domain is a singleton ($|D_u| = 1$), the colour is forced and the display uses
the wording ``must be''.  Conditional entries --- colours available only under specific human
assignments --- are flagged separately:
\begin{center}
\small
\texttt{Possible: a5 = Green OR Red}\\
\texttt{Conditional: a5 = Blue IF h4 $\neq$ Blue}
\end{center}

\paragraph{Design intent.}
This condition highlights remaining degrees of freedom for agent nodes, making it straightforward
to see how much flexibility the system retains.  It does not, however, directly express the
causal link from human decisions to agent constraints.

\subsection{Condition~3: User-Centric Natural Language}
\label{sec:c3}

The same user-centric consequence information as Condition~1 is presented, but rendered into
natural language by the LLM rather than as a logical formula.

\paragraph{Display format.}
Human nodes display a plain-English summary:
\begin{center}
\small
\texttt{Because h4 is Blue, the neighbouring agent nodes are restricted}\\
\texttt{to two configurations: either a5 is Red and a4 is Green,}\\
\texttt{or a5 is Green and a4 is Red.}
\end{center}
When few constraints are active, the summary may convey that the agent remains largely
unconstrained:
\begin{center}
\small
\texttt{Your current choice leaves most agent configurations open.}\\
\texttt{The neighbouring nodes are only weakly constrained.}
\end{center}

\paragraph{Design intent.}
This condition tests whether LLM-rendered descriptions reduce the cognitive load of interpreting
formulaic logical expressions while preserving the user-centric causal framing.

\subsection{Condition~4: Agent-Centric Natural Language}
\label{sec:c4}

The same agent-centric domain information as Condition~2 is rendered into natural language by
the LLM.

\paragraph{Display format.}
Agent nodes display plain-English descriptions of their feasible colours:
\begin{center}
\small
\texttt{a5 can currently be Red or Green. Blue would only become}\\
\texttt{available if h4 changed to a different colour.}
\end{center}

\paragraph{Design intent.}
This condition tests whether natural-language framing of remaining options is sufficient for users
to guide their colour choices efficiently, without the additional cognitive overhead of interpreting
consequence-directed formulae.

%% ============================================================
\section{Experimental Design}
\label{sec:experiment}
%% ============================================================

\subsection{Task}

Participants must colour all nodes in their cluster such that a valid global colouring exists ---
that is, such that $S_{a_1}(\alpha_h) \neq \emptyset$ and $S_{a_2}(\alpha_h) \neq \emptyset$
simultaneously.  Participants may:
\begin{itemize}
  \item Change any node colour freely at any time.
  \item Leave nodes grey (unassigned) while exploring possibilities.
  \item Consult the constraint panels at any time.
\end{itemize}
The task ends when the participant confirms that a feasible assignment exists and submits their
final colouring.  There is no dialogue, no agent messaging, and no turn structure.

\subsection{Conditions and Design}

We use a within-subjects design: each participant completes the task under all four conditions,
counterbalanced across participants using a $4 \times 4$ Latin square to control for learning and
order effects.  A different graph instance is used for each condition to prevent transfer of
specific solution knowledge.

\subsection{Graph Instances}

Four distinct graph instances are generated, matched for:
\begin{itemize}
  \item $|V| = 15$, five nodes per cluster.
  \item $|\mathcal{C}| = 3$ colours.
  \item 4--6 inter-cluster boundary edges.
  \item Chromatic number equal to 3, ensuring a valid 3-colouring exists but requiring
        coordination between clusters.
  \item Comparable solution density: number of valid colourings within 10\% across instances.
\end{itemize}
Graph instances are generated by random sampling with these constraints and validated by
exhaustive enumeration.

\subsection{Procedure}

The session proceeds as follows:
\begin{enumerate}
  \item \textbf{Introduction} (10 min): task explanation, graph colouring rules, interface tutorial
        using a practice graph not used in the main conditions.
  \item \textbf{Condition blocks} (4 $\times$ 15 min): one block per condition, separated by
        2-minute breaks.  Each block includes:
        \begin{enumerate}[label=\alph*.]
          \item Brief instruction screen showing the current condition's interface layout.
          \item Task on the assigned graph instance.
          \item Post-block questionnaire (workload, clarity, control).
        \end{enumerate}
  \item \textbf{End-of-session battery} (15 min): full NASA-TLX, trust questionnaire, preference
        ranking across all four conditions.
\end{enumerate}
Total session duration is approximately 90 minutes.

\subsection{Participants}

Approximately 32 participants ($n \geq 28$ after exclusions), aged 18+, with no prior experience
with graph colouring problems.  Participants are recruited from a university participant pool.
Inclusion criteria: normal or corrected-to-normal colour vision (essential for graph display),
native or near-native English (for natural-language conditions), and ability to complete a
90-minute session.

\subsection{Measures}

\paragraph{Objective.}
\begin{itemize}
  \item \textbf{Time to feasibility}: wall-clock time from task start to confirmed feasible
        submission.
  \item \textbf{Node adjustments}: total number of colour changes made during the task.
  \item \textbf{Infeasible visits}: number of timesteps during which the global configuration
        was infeasible ($S_{a_1} = \emptyset$ or $S_{a_2} = \emptyset$).
  \item \textbf{Exploration breadth}: number of distinct colour assignments visited across
        human nodes, as a proxy for search strategy.
\end{itemize}

\paragraph{Subjective (post-block).}
\begin{itemize}
  \item \textbf{Workload}: NASA-TLX~\citep{hart1988development} six subscales, administered after
        each block.
  \item \textbf{Perceived clarity}: custom 5-item scale measuring how clearly the interface
        communicated the constraints.
  \item \textbf{Perceived control}: adapted from~\citep{hoffman2019evaluating}; degree to which
        the participant felt in control of the outcome.
\end{itemize}

\paragraph{Subjective (end of session).}
\begin{itemize}
  \item \textbf{Trust in system guidance}: post-task trust questionnaire adapted
        from~\citep{jian2000foundations}, assessing trust in the constraint information displayed.
  \item \textbf{Mode preference}: paired-comparison ranking across all four conditions, eliciting
        which interface the participant would prefer for sustained use.
  \item \textbf{Open comments}: optional free-text field for qualitative feedback.
\end{itemize}

%% ============================================================
\section{Expected Results and Hypotheses}
\label{sec:results}
%% ============================================================

The $2 \times 2$ design permits main effects and interactions between perspective and representation.
We state the primary hypotheses below.

\begin{description}[leftmargin=2em]
  \item[H1 \emph{(Representation -- Workload)}]
    LLM-generated natural-language summaries (Conditions 3 and 4) will yield lower NASA-TLX
    scores than formulaic conditions (Conditions 1 and 2), because they reduce the syntactic
    parsing overhead of logical expressions.
  \item[H2 \emph{(Perspective -- Convergence)}]
    Agent-centric conditions (Conditions 2 and 4) will yield faster time to feasibility, because
    directly displaying remaining degrees of freedom gives the human more immediately actionable
    information about which node changes are most effective.
  \item[H3 \emph{(Perspective -- Understanding)}]
    User-centric conditions (Conditions 1 and 3) will yield higher perceived clarity scores,
    because linking the human's choices directly to agent outcomes makes the causal structure of
    the constraint system more legible.
  \item[H4 \emph{(Interaction)}]
    The benefit of natural-language rendering will be larger for user-centric conditions, because
    the consequence descriptions produced in user-centric mode are structurally more complex
    (disjunctions over joint agent assignments) and thus benefit more from summarisation.
  \item[H5 \emph{(Trust)}]
    Trust in system guidance will be higher for formulaic conditions, as the precise logical
    statements are more verifiable than LLM-generated prose, which participants may perceive as
    potentially imprecise.
\end{description}

Table~\ref{tab:expected} summarises expected qualitative performance across the four conditions.

\begin{table}[t]
\centering
\caption{Expected qualitative performance of the four conditions.
VH = Very High, H = High, M = Medium, L = Low, VL = Very Low.}
\label{tab:expected}
\begin{tabular}{lcccc}
\toprule
\textbf{Attribute} & \textbf{C1 (UC-F)} & \textbf{C2 (AC-F)} & \textbf{C3 (UC-NL)}
                   & \textbf{C4 (AC-NL)} \\
\midrule
Time to feasibility (low = fast) & H   & L   & M   & VL  \\
Node adjustments                 & H   & M   & M   & L   \\
Infeasible visits                & M   & L   & M   & VL  \\
Cognitive workload               & VH  & H   & M   & L   \\
Perceived clarity                & H   & M   & VH  & M   \\
Perceived control                & M   & H   & M   & H   \\
Trust in guidance                & H   & H   & M   & M   \\
Mode preference                  & L   & M   & M   & VH  \\
\bottomrule
\end{tabular}
\end{table}

\noindent (UC = User-Centric; AC = Agent-Centric; F = Formulaic; NL = Natural Language.)

We expect Condition~4 (agent-centric, natural-language) to offer the best overall trade-off
between speed and usability, and Condition~1 (user-centric, formulaic) to provide the highest
precision at the cost of high cognitive load.  The primary scientific question is whether
user-centric or agent-centric framing more strongly improves human understanding when natural
language is available.

%% ============================================================
\section{Conclusion}
\label{sec:conclusion}
%% ============================================================

This study shifts the focus of human--agent coordination research from dialogue-based negotiation
to constraint explanation.  Rather than asking how agents and humans should exchange proposals, we
ask how AI systems should present the constraint structure of a problem so that humans can make
effective decisions.  The redesigned platform isolates two dimensions of this question ---
perspective and representation --- and provides an experimental infrastructure for testing their
effects on understanding, workload, and performance.

The interface designed here --- interactive constraint visualisation with real-time feasibility
computation --- represents a novel class of human--AI interaction that could generalise beyond
graph colouring to any cooperative constraint-satisfaction task where a human makes partial
assignments and agents must adapt.

Future work includes extending the approach to larger DCOP instances where exhaustive enumeration
must be replaced by approximate projection, studying how participants develop constraint-reasoning
strategies over repeated trials, and investigating richer explanation modalities such as
counterfactual comparisons (``if you had chosen green instead of blue\ldots'').

\bibliographystyle{plainnat}
\bibliography{refs}

%% ============================================================
%% TECHNICAL APPENDIX
%% (Detailed formal definitions, algorithms, and implementation notes)
%% ============================================================

\clearpage
\appendix
\renewcommand{\thesection}{\Alph{section}}

\begin{center}
  {\LARGE\bfseries Technical Appendix}\\[0.5em]
  {\large Formal Definitions, Algorithms, and Implementation Details}
\end{center}
\vspace{1em}

\noindent This appendix provides complete formal definitions, pseudocode for all algorithms,
implementation details, and example interface traces.  It is intended as a self-contained
technical reference for the system described in the main paper.

\tableofcontents
\newpage

%% ============================================================
\section{Formal Problem Definition}
\label{app:problem}
%% ============================================================

\subsection{Graph Colouring DCOP}

\begin{definition}[Graph Colouring DCOP]
A graph colouring DCOP is a tuple $\langle G, \mathcal{C}, \mathcal{P}, \mathcal{V} \rangle$:
\begin{itemize}
  \item $G = (V, E)$: undirected graph with node set $V$, edge set $E$.
  \item $\mathcal{C}$: finite colour domain, $|\mathcal{C}| = k$ (default $k=3$).
  \item $\mathcal{P} = \{a_1, a_2, h\}$: participants (agents $a_1, a_2$, human $h$).
  \item $\mathcal{V}: \mathcal{P} \to 2^V$: cluster partition,
        $\bigsqcup_{p \in \mathcal{P}} \mathcal{V}(p) = V$.
\end{itemize}
\end{definition}

\begin{definition}[Boundary and Internal Nodes]
For each participant $p$ with cluster $V_p$:
\begin{align*}
  \partial V_p &= \{v \in V \setminus V_p \mid \exists\, u \in V_p: (u,v) \in E\}
  &&\text{(visible external nodes)}\\
  B_p &= \{u \in V_p \mid \exists\, v \in \partial V_p: (u,v) \in E\}
  &&\text{(boundary nodes of }p\text{)}\\
  I_p &= V_p \setminus B_p
  &&\text{(internal nodes of }p\text{)}
\end{align*}
\end{definition}

\begin{definition}[Partial Human Assignment]
The human's configuration is a function
$\alpha_h: V_h \to \mathcal{C} \cup \{\bot\}$,
where $\bot$ denotes an unassigned (grey) node.
The effective domain is $\mathrm{dom}(\alpha_h) = \{v \in V_h \mid \alpha_h(v) \neq \bot\}$.
The \emph{colour projection} of $\alpha_h$ onto a set $U \subseteq V_h$ is
$\alpha_h|_U = \{v \mapsto \alpha_h(v) \mid v \in U \cap \mathrm{dom}(\alpha_h)\}$.
\end{definition}

\begin{definition}[Agent Feasibility Set]
For agent $a$ with cluster $V_a$ and boundary $\partial V_a \cap V_h \subseteq V_h$,
the feasibility set under partial human assignment $\alpha_h$ is:
\[
  S_a(\alpha_h) = \bigl\{ x_a \in \mathcal{C}^{|V_a|}
    \;\big|\; J_a(x_a,\, \alpha_h|_{\partial V_a}) = 0 \bigr\}
\]
where $J_a(x_a, \alpha_h|_{\partial V_a})$ counts conflicts on edges incident to $V_a$, treating
unassigned boundary neighbours as imposing no constraint.
\end{definition}

\begin{definition}[Domain Projection]
For agent node $u \in V_a$, the \emph{domain projection} under $\alpha_h$ is:
\[
  D_u(\alpha_h) = \bigl\{ c \in \mathcal{C} \;\big|\;
    \exists\, x_a \in S_a(\alpha_h): x_a(u) = c \bigr\}
\]
\end{definition}

\begin{definition}[Consequence Set]
For human node $v \in \mathrm{dom}(\alpha_h)$ with $\alpha_h(v) = c$, the \emph{consequence set}
for agent $a$ is:
\[
  \mathrm{Cons}_{a,v}(\alpha_h) =
    \bigl\{ x_a \in S_a(\alpha_h) \;\big|\; \text{consistent with } \alpha_h(v) = c \bigr\}
\]
In practice this is simply $S_a(\alpha_h)$ restricted to configurations that do not conflict with
the fixed colour of $v$'s boundary neighbours in $V_a$.
\end{definition}

\subsection{Notation Summary}

\begin{itemize}[leftmargin=*]
  \item $\alpha_h$: partial human assignment; $\alpha_h(v) = \bot$ if $v$ is grey.
  \item $\mathrm{dom}(\alpha_h)$: fixed (non-grey) human nodes.
  \item $S_a(\alpha_h)$: feasibility set for agent cluster $a$.
  \item $D_u(\alpha_h)$: feasible colour domain for agent node $u$.
  \item $J(x) = \sum_{(u,v)\in E} \mathbf{1}[x_u=x_v]$: global conflict penalty.
  \item $J_a(x_a, \alpha_h|_{\partial V_a})$: local penalty for agent $a$'s edges.
  \item $k = |\mathcal{C}|$: number of colours (default 3).
  \item $n = |V_a|$: nodes per cluster (default 5).
  \item $\epsilon = 10^{-6}$: numerical tolerance for zero-penalty checks.
\end{itemize}

%% ============================================================
\section{Real-Time Constraint Computation}
\label{app:computation}
%% ============================================================

\subsection{Feasibility Set Enumeration}

\begin{algorithm}
\caption{ComputeFeasibilitySet: exact enumeration of $S_a(\alpha_h)$}
\label{alg:feasset}
\begin{algorithmic}[1]
\REQUIRE Agent $a$, partial assignment $\alpha_h$, graph $(V, E)$, colour domain $\mathcal{C}$
\STATE $\mathrm{bnd} \gets \partial V_a \cap \mathrm{dom}(\alpha_h)$
  \COMMENT{boundary nodes with fixed human colours}
\STATE $S \gets \emptyset$
\FORALL{$(c_1, \ldots, c_n) \in \mathcal{C}^{|V_a|}$ \COMMENT{$3^5 = 243$ candidates}}
  \STATE $x_a \gets \{v_i \mapsto c_i\}_{i=1}^{n}$
  \STATE $\mathrm{conflicts} \gets 0$
  \FORALL{$(u,v) \in E$ with $u \in V_a$}
    \IF{$v \in \mathrm{bnd}$ \AND $x_a[u] = \alpha_h(v)$}
      \STATE $\mathrm{conflicts} \gets \mathrm{conflicts} + 1$
    \ENDIF
    \IF{$v \in V_a$ \AND $x_a[u] = x_a[v]$}
      \STATE $\mathrm{conflicts} \gets \mathrm{conflicts} + 1$
    \ENDIF
  \ENDFOR
  \IF{$\mathrm{conflicts} = 0$}
    \STATE $S \gets S \cup \{x_a\}$
  \ENDIF
\ENDFOR
\RETURN $S$
\end{algorithmic}
\end{algorithm}

\subsection{Domain Projection}

\begin{algorithm}
\caption{ComputeDomainProjection: $D_u(\alpha_h)$ for all agent nodes}
\label{alg:domain}
\begin{algorithmic}[1]
\REQUIRE Agent $a$, feasibility set $S_a(\alpha_h)$
\STATE Initialise $D_u \gets \emptyset$ for all $u \in V_a$
\FORALL{$x_a \in S_a(\alpha_h)$}
  \FORALL{$u \in V_a$}
    \STATE $D_u \gets D_u \cup \{x_a[u]\}$
  \ENDFOR
\ENDFOR
\RETURN $\{u \mapsto D_u \mid u \in V_a\}$
\end{algorithmic}
\end{algorithm}

\subsection{Conditional Domain Entries}

A colour $c \in \mathcal{C} \setminus D_u(\alpha_h)$ is \emph{conditionally available} for node
$u$ if there exists an alternative partial human assignment $\alpha_h'$ that differs from
$\alpha_h$ in at most one node and makes $c \in D_u(\alpha_h')$.  This is used in
Condition~2 to populate the conditional entries of the domain display.

\begin{algorithm}
\caption{ComputeConditionalDomains: single-change conditional entries}
\label{alg:conditional}
\begin{algorithmic}[1]
\REQUIRE Agent $a$, partial assignment $\alpha_h$, domain $D_u(\alpha_h)$
\STATE $\mathrm{Cond} \gets \{u \mapsto \{\} \mid u \in V_a\}$
\FORALL{$v \in \mathrm{dom}(\alpha_h)$ with $v \in \partial V_a$}
  \FORALL{$c' \in \mathcal{C} \setminus \{\alpha_h(v)\}$}
    \STATE $\alpha_h' \gets \alpha_h[v \mapsto c']$
    \STATE $S' \gets \text{ComputeFeasibilitySet}(a, \alpha_h')$
    \STATE $D' \gets \text{ComputeDomainProjection}(a, S')$
    \FORALL{$u \in V_a$}
      \FORALL{$d \in D'[u] \setminus D_u(\alpha_h)$}
        \STATE $\mathrm{Cond}[u] \gets \mathrm{Cond}[u] \cup \{(d,\, v \neq \alpha_h(v))\}$
      \ENDFOR
    \ENDFOR
  \ENDFOR
\ENDFOR
\RETURN $\mathrm{Cond}$
\end{algorithmic}
\end{algorithm}

\subsection{Real-Time Update Pipeline}

\begin{algorithm}
\caption{OnNodeColourChange: pipeline triggered by human interaction}
\label{alg:pipeline}
\begin{algorithmic}[1]
\REQUIRE Updated node $v$, new colour $c'$ (or $\bot$ if reset to grey), current $\alpha_h$
\STATE $\alpha_h[v] \gets c'$
\FORALL{$a \in \{a_1, a_2\}$ (in parallel)}
  \STATE $S_a \gets \text{ComputeFeasibilitySet}(a, \alpha_h)$
  \STATE $D_a \gets \text{ComputeDomainProjection}(a, S_a)$
  \IF{active condition is agent-centric}
    \STATE $\mathrm{Cond}_a \gets \text{ComputeConditionalDomains}(a, \alpha_h, D_a)$
  \ELSE
    \STATE $\mathrm{Cons}_a \gets \text{ComputeConsequenceSets}(a, \alpha_h, S_a)$
  \ENDIF
  \STATE $\mathrm{Feasible}_a \gets (S_a \neq \emptyset)$
\ENDFOR
\STATE Update feasibility panel with $\mathrm{Feasible}_{a_1}$, $\mathrm{Feasible}_{a_2}$
\STATE Update node annotations according to active condition
\IF{active condition uses LLM representation}
  \STATE Queue LLM summary calls for changed annotations (Section~\ref{app:llmdetail})
\ENDIF
\end{algorithmic}
\end{algorithm}

%% ============================================================
\section{Infeasibility Repair}
\label{app:repair}
%% ============================================================

\begin{algorithm}
\caption{InfeasibilityRepair: minimal node set to restore feasibility}
\label{alg:repair}
\begin{algorithmic}[1]
\REQUIRE Agent $a$, partial assignment $\alpha_h$
\STATE $\mathrm{fixed} \gets \mathrm{dom}(\alpha_h) \cap \partial V_a$
  \COMMENT{human nodes adjacent to agent $a$}
\FOR{$r = 1, 2, \ldots, |\mathrm{fixed}|$}
  \FORALL{$R \subseteq \mathrm{fixed}$ with $|R| = r$}
    \STATE $\alpha_h' \gets \alpha_h|_{V_h \setminus R}$
      \COMMENT{release $R$ to unassigned}
    \IF{$\text{ComputeFeasibilitySet}(a, \alpha_h') \neq \emptyset$}
      \STATE \textbf{return} $R$ with suggested colours for each node in $R$
    \ENDIF
  \ENDFOR
\ENDFOR
\RETURN ``No valid colouring exists for this graph.''
  \COMMENT{only if graph is not 3-colourable; guaranteed not to happen by construction}
\end{algorithmic}
\end{algorithm}

The suggested colours for each node $v \in R$ are drawn from a feasible configuration of the
smallest $S_a(\alpha_h')$ found, projected onto the boundary neighbours of $v$.

%% ============================================================
\section{LLM Summary Generation}
\label{app:llmdetail}
%% ============================================================

\subsection{Input Construction}

For each node annotation that changes in an LLM-representation condition, the system constructs
a structured prompt as follows.

\paragraph{Agent-centric (Condition~4).}
The prompt provides:
\begin{lstlisting}
{
  "node": "a5",
  "feasible_colours": ["Red", "Green"],
  "conditional_colours": [
    {"colour": "Blue", "condition": "h4 != Blue"}
  ],
  "human_assignment": {"h4": "Blue", "h3": "Red"}
}
\end{lstlisting}
System instruction: \emph{``Describe, in one or two plain-English sentences, what colours node
a5 can take and under what conditions. Do not introduce constraints not present in the data.''}

\paragraph{User-centric (Condition~3).}
The prompt provides the consequence set for the changed human node:
\begin{lstlisting}
{
  "human_node": "h4",
  "human_colour": "Blue",
  "agent": "a1",
  "feasible_agent_configs": [
    {"a4": "Green", "a5": "Red"},
    {"a4": "Red",   "a5": "Green"}
  ]
}
\end{lstlisting}
System instruction: \emph{``Summarise, in one or two plain-English sentences, what agent
configurations are implied by this human node choice. Do not introduce constraints not present
in the data.''}

\subsection{LLM Call Specification}

All LLM calls use the chat completions API with:
\begin{itemize}
  \item \texttt{temperature = 0.2} (low variance for reproducibility).
  \item \texttt{max\_tokens = 80} (enforces brevity).
  \item Model: configurable; default is a capable instruction-following model.
\end{itemize}
The LLM response is used as-is.  If the API call fails, the system falls back to the formulaic
display for that node, and logs a warning.  This is the only fallback permitted, and is
justified because the fallback is a more informative display, not a degraded one.

\subsection{Hallucination Logging}

All LLM inputs and outputs are logged to \texttt{results/llm\_explanations.jsonl}, including:
\begin{itemize}
  \item The exact structured input (feasibility set or domain projection).
  \item The LLM prompt and response.
  \item A post-hoc verification flag set by comparing the LLM output against the structured
        input for any mentioned constraint that is absent from the input.
\end{itemize}
Post-hoc analysis of this log enables detection of hallucinated constraints that were presented
to participants during the study.

%% ============================================================
\section{Interface Condition Specifications}
\label{app:conditions}
%% ============================================================

\subsection{Condition~1: User-Centric Formulaic}

Each human node $v \in \mathrm{dom}(\alpha_h)$ with $\alpha_h(v) = c$ displays:

\noindent\textbf{If $|\mathrm{Cons}_{a,v}(\alpha_h)| = 1$} (forced assignment):
\begin{lstlisting}
Because: h4 = Blue
Forces: a5 = Red AND a4 = Green
\end{lstlisting}

\noindent\textbf{If $|\mathrm{Cons}_{a,v}(\alpha_h)| > 1$} (disjunctive):
\begin{lstlisting}
Because: h4 = Blue
Allows: (a5 = Red AND a4 = Green)
        OR (a5 = Green AND a4 = Red)
[+ 3 more configurations]
\end{lstlisting}
When more than 4 configurations exist, the display shows the first 4 and a count of remaining ones.

\noindent\textbf{If $|\mathrm{Cons}_{a,v}(\alpha_h)| = 0$} (infeasible):
\begin{lstlisting}
Because: h4 = Blue
INFEASIBLE: No valid agent configuration exists.
[Change suggestion displayed below]
\end{lstlisting}

\subsection{Condition~2: Agent-Centric Formulaic}

Each agent node $u \in V_a$ displays:
\begin{lstlisting}
a5
Possible:     Green, Red
Conditional:  Blue IF h4 != Blue
\end{lstlisting}
When $|D_u(\alpha_h)| = 0$, the node annotation displays \texttt{[INFEASIBLE]} in red.

\subsection{Condition~3: User-Centric Natural Language}

Human node annotations are replaced by LLM-generated summaries of the consequence sets
(Section~\ref{app:llmdetail}).  Layout and positioning are identical to Condition~1.

\subsection{Condition~4: Agent-Centric Natural Language}

Agent node annotations are replaced by LLM-generated summaries of the domain projections
(Section~\ref{app:llmdetail}).  Layout and positioning are identical to Condition~2.

%% ============================================================
\section{Example Interaction Traces}
\label{app:traces}
%% ============================================================

The following traces illustrate the system's behaviour across all four conditions for an identical
sequence of human interactions on a sample graph.

\subsection{Scenario}

Nodes $h_3, h_4, h_5$ are in the human cluster; $a_4, a_5$ are boundary nodes of agent $a_1$.
Boundary edge $(h_4, a_5)$ exists.  The human begins with all nodes grey.

\paragraph{Step 1: Human sets $h_4 = \textit{Blue}$.}

\noindent \textbf{Condition~1 output (user-centric formulaic):}
\begin{lstlisting}
h4 [Blue]
Because: h4 = Blue
Allows: (a5 = Red AND a4 = Green)
        OR (a5 = Green AND a4 = Red)
\end{lstlisting}

\noindent \textbf{Condition~2 output (agent-centric formulaic):}
\begin{lstlisting}
a5: Possible: {Red, Green}   Conditional: Blue IF h4 != Blue
a4: Possible: {Red, Green, Blue}
\end{lstlisting}

\noindent \textbf{Condition~3 output (user-centric natural language):}
\begin{lstlisting}
h4 [Blue]
Setting h4 to Blue means the neighbouring agent must choose
between two configurations: a5 Red with a4 Green, or a5 Green
with a4 Red.
\end{lstlisting}

\noindent \textbf{Condition~4 output (agent-centric natural language):}
\begin{lstlisting}
a5: Can currently be Red or Green. Blue would only be possible
    if h4 were changed to a different colour.
a4: Can be any colour.
\end{lstlisting}

\paragraph{Step 2: Human sets $h_4 = \textit{Red}$ (changes colour).}
Agent feasibility set updates; all conditions update immediately.  In Condition~2, $a_5$'s domain
now excludes Red and adds Blue as a directly possible colour.

\paragraph{Step 3: Human sets $h_3 = \textit{Blue}$, creating an infeasible state.}
Global feasibility panel updates to show \textsc{Infeasible} for $a_1$.  In all conditions, the
repair proposal is displayed:
\begin{lstlisting}
INFEASIBLE: No valid colouring for Agent 1.
Suggested fix: change h3 to Green or Red.
\end{lstlisting}

%% ============================================================
\section{Graph Instance Generation}
\label{app:graphs}
%% ============================================================

The four graph instances used in the study are generated by the following procedure:

\begin{algorithm}
\caption{GenerateMatchedGraphInstances}
\label{alg:graphgen}
\begin{algorithmic}[1]
\REQUIRE Number of instances $M=4$, nodes per cluster $n=5$, colours $k=3$
\STATE $\mathrm{instances} \gets \emptyset$
\WHILE{$|\mathrm{instances}| < M$}
  \STATE Sample random intra-cluster edges for each cluster (Erd\H{o}s--R\'{e}nyi, $p=0.4$)
  \STATE Sample 4--6 inter-cluster boundary edges between $B_h$ and $\partial V_h$
  \STATE Verify $G$ is $k$-colourable by exhaustive search
  \STATE Count the number of valid $k$-colourings $N_{\mathrm{valid}}$
  \IF{$N_{\mathrm{valid}} \geq N_{\min}$ \AND $N_{\mathrm{valid}} \leq N_{\max}$}
    \STATE $\mathrm{instances} \gets \mathrm{instances} \cup \{G\}$
  \ENDIF
\ENDWHILE
\STATE Verify that valid-colouring counts across all instances are within 10\% of each other
\RETURN $\mathrm{instances}$
\end{algorithmic}
\end{algorithm}

The bounds $N_{\min}$ and $N_{\max}$ are set to ensure instances are neither trivially easy
(many valid colourings) nor excessively constrained (very few).  Exact values are determined
during piloting.

\end{document}
